\documentclass{article}
\usepackage[utf8]{inputenc}
\usepackage[T1]{fontenc}
\usepackage[spanish]{babel}
\usepackage{amsmath, amssymb}
\usepackage{enumitem}
\usepackage{geometry}
\geometry{margin=2.5cm}

\title{Entrega Semanal 3 --- Esbozo de Preliminares} \author{Noel Pérez Calvo} \date{\today}

\begin{document}

\maketitle

\section*{Preliminares: qué se va a decir de cada elemento}

\begin{itemize}[leftmargin=1.5em]

  \item Ecuaciones no lineales
  \begin{itemize}
\item definición de ecuación no lineal (ecuación donde la incógnita aparece con grado mayor que 1 o dentro de funciones no lineales) \item diferencia con ecuaciones lineales \item ecuación paramétrica: $F(x; \alpha_1, \ldots, \alpha_n) = 0$, donde $x$ es la incógnita y los $\alpha_i$ son parámetros \item el concepto de raíz o cero de una ecuación \item que una ecuación puede tener múltiples raíces reales
  \end{itemize}

  \item Expresiones analíticas
  \begin{itemize}
\item qué es una expresión analítica (fórmula cerrada) \item diferencia entre solución numérica y solución analítica \item ejemplo: la fórmula general cuadrática $x = \frac{-b \pm \sqrt{b^2 - 4ac}}{2a}$ como expresión analítica de las raíces de $ax^2+bx+c=0$ \item teorema de Abel--Ruffini: no existe fórmula general para polinomios de grado $\geq 5$, lo cual motiva un enfoque computacional
  \end{itemize}

  \item Resolución numérica
  \begin{itemize}
\item qué es resolver numéricamente una ecuación (obtener un valor decimal aproximado) \item concepto de método iterativo para encontrar raíces (Newton--Raphson como ejemplo clásico) \item necesidad de un punto inicial (initial guess) y posibilidad de no convergencia \item que en este trabajo se usa como paso intermedio: primero se resuelve numéricamente para muchos parámetros y luego se descubre la fórmula
  \end{itemize}

  \item Rama de soluciones
  \begin{itemize}
\item cuando una ecuación tiene $k$ raíces, se ordenan de menor a mayor \item la $i$-ésima raíz más pequeña para cada combinación de parámetros forma la rama $i$ \item cada rama se trata como un problema de regresión independiente
  \end{itemize}

  \item SymPy
  \begin{itemize}
\item biblioteca de Python para matemática simbólica \item permite representar y manipular expresiones algebraicas de forma exacta \item funciones \texttt{solve} (solución simbólica exacta) y \texttt{nsolve} (solución numérica) \item función \texttt{sympify} para convertir texto a expresiones simbólicas
  \end{itemize}

  \item Parser simbólico
  \begin{itemize}
\item programa que convierte una cadena de texto en una estructura simbólica manipulable \item en el contexto del trabajo: recibe la ecuación como string y produce un objeto SymPy \item detecta automáticamente variables e incógnitas
  \end{itemize}

  \item Grid de parámetros y producto cartesiano
  \begin{itemize}
\item discretización de un rango continuo en puntos equiespaciados \item producto cartesiano: todas las combinaciones posibles de valores de cada parámetro \item ejemplo: si $a \in \{1,2,3\}$ y $b \in \{4,5\}$, el producto cartesiano es $\{(1,4),(1,5),(2,4),(2,5),(3,4),(3,5)\}$
  \end{itemize}

  \item Regresión simbólica
  \begin{itemize}
\item técnica de aprendizaje automático que busca la expresión matemática que mejor se ajusta a un conjunto de datos \item a diferencia de la regresión clásica (que ajusta coeficientes de un modelo fijo), la regresión simbólica descubre la estructura de la fórmula \item las expresiones se representan como árboles de expresiones
  \end{itemize}

  \item Programación genética
  \begin{itemize}
\item metaheurística de optimización inspirada en la evolución biológica \item los individuos son programas (en este caso, árboles de expresiones matemáticas) \item operaciones: mutación (modificar un nodo del árbol), crossover (intercambiar subárboles entre dos individuos), selección \item es el mecanismo que usa PySR internamente para buscar expresiones
  \end{itemize}

  \item Población y ciclo evolutivo
  \begin{itemize}
\item población: conjunto de individuos (expresiones candidatas) que evolucionan simultáneamente \item ciclo evolutivo: una iteración de selección, crossover, mutación y reemplazo \item modelo de islas: múltiples poblaciones independientes con migración periódica entre ellas para mantener diversidad
  \end{itemize}

  \item Tournament selection y presión selectiva
  \begin{itemize}
\item tournament selection: se eligen $k$ individuos al azar y el mejor pasa a la siguiente generación \item presión selectiva: qué tan agresivamente se favorece a los mejores individuos \item mayor tamaño de torneo $\Rightarrow$ mayor presión $\Rightarrow$ convergencia más rápida pero menor diversidad \item en el trabajo se redujo la presión selectiva para evitar convergencia prematura
  \end{itemize}

  \item Diversidad de la población y convergencia prematura
  \begin{itemize}
\item diversidad: variedad genética (estructural) entre los individuos de la población \item convergencia prematura: cuando la población pierde diversidad demasiado rápido y se queda atrapada en un óptimo local \item parámetros que afectan la diversidad: tamaño de torneo, tasa de migración, probabilidad de crossover
  \end{itemize}

  \item Crossover
  \begin{itemize}
\item operador genético que combina material de dos individuos padres para crear descendencia \item en programación genética: intercambio de subárboles entre dos árboles de expresiones \item la tasa de crossover controla el balance entre exploración y explotación
  \end{itemize}

  \item Exploración vs.\ explotación
  \begin{itemize}
\item exploración: buscar en regiones nuevas del espacio de soluciones \item explotación: refinar soluciones ya encontradas \item un buen algoritmo evolutivo necesita equilibrio entre ambas
  \end{itemize}

  \item Hall of fame y frente de Pareto
  \begin{itemize}
\item hall of fame en PySR: conjunto de las mejores expresiones encontradas a lo largo de la búsqueda \item frente de Pareto: conjunto de soluciones no dominadas en el compromiso complejidad vs.\ precisión \item una expresión está en el frente si no hay otra que sea simultáneamente más simple y más precisa
  \end{itemize}

  \item Ecuaciones candidatas
  \begin{itemize}
\item cada expresión en el hall of fame es una ecuación candidata \item en el trabajo se evalúan todas las candidatas (no solo la ``mejor'') porque la que PySR considera mejor por balance complejidad/error no siempre es la que cubre más puntos
  \end{itemize}

  \item Operadores matemáticos seguros (\texttt{safe\_pow})
  \begin{itemize}
\item problema: durante la evolución genética, expresiones intermedias pueden producir NaN (raíz cuadrada de negativo, por ejemplo) \item solución: definir $\texttt{safe\_pow}(x, y) = \text{sign}(x) \cdot |x|^y$ \item permite descubrir raíces de cualquier orden sin errores numéricos \item PySR tiene un mecanismo llamado bumper para esto, pero en la práctica eliminaba expresiones con raíces antes de que pudieran componerse
  \end{itemize}

  \item MSE (error cuadrático medio)
  \begin{itemize}
\item función de pérdida: $\text{MSE} = \frac{1}{n}\sum_{i=1}^{n}(y_i - \hat{y}_i)^2$ \item mide la diferencia promedio al cuadrado entre valores reales y predichos \item se eligió MSE estándar porque los datos generados numéricamente no tienen ruido
  \end{itemize}

  \item Ruido en datos
  \begin{itemize}
\item datos con ruido: contienen errores aleatorios (propios de mediciones experimentales) \item datos sin ruido: valores exactos (como los generados numéricamente en este trabajo) \item la ausencia de ruido permite usar MSE sin regularización adicional
  \end{itemize}

  \item Batching (mini-batches)
  \begin{itemize}
\item técnica de dividir el dataset en subconjuntos pequeños para entrenar \item reduce el consumo de memoria RAM \item en el trabajo se usan mini-batches de 200 puntos
  \end{itemize}

  \item Expresión piecewise
  \begin{itemize}
\item función definida por partes: distintas fórmulas aplican en distintas regiones del dominio \item surge cuando la regresión simbólica iterativa encuentra varias expresiones para una misma rama, cada una cubriendo una región distinta del espacio de parámetros \item se ensamblan automáticamente detectando las fronteras de cada región
  \end{itemize}

  \item Benchmark
  \begin{itemize}
\item conjunto de problemas de prueba estandarizados para evaluar un sistema \item en el trabajo: 43 casos de prueba (10 lineales, 14 cuadráticos, 6 cúbicos, 4 cuárticos, 1 quíntico, 8 especiales) \item cada caso tiene: la ecuación, las raíces esperadas y métricas automáticas \item relación con SRBench (benchmark estándar de regresión simbólica, NeurIPS 2022)
  \end{itemize}

  \item Bases de Gröbner
  \begin{itemize}
\item herramienta del álgebra computacional para manipular ideales polinómicos \item análogas a la eliminación gaussiana, pero para sistemas de ecuaciones polinómicas no lineales \item permiten decidir si un polinomio pertenece a un ideal dado (problema de pertenencia) \item se calculan mediante el algoritmo de Buchberger \item en el trabajo: se usan para verificar algebraicamente que las expresiones descubiertas son raíces exactas de la ecuación original
  \end{itemize}

  \item Bibliotecas de regresión simbólica (PySR, gplearn)
  \begin{itemize}
\item PySR: biblioteca basada en Julia con interfaz Python, usa modelo de islas y exploración multi-población \item gplearn: biblioteca de Python pura, basada en scikit-learn \item se seleccionó PySR por su rendimiento superior en benchmarks y su backend compilado en Julia
  \end{itemize}

  \item Julia
  \begin{itemize}
\item lenguaje de programación de alto rendimiento para computación científica \item PySR lo usa como backend para que la búsqueda evolutiva sea eficiente \item compilación JIT (just-in-time) permite velocidad comparable a C
  \end{itemize}

  \item Estado del arte
  \begin{itemize}
\item revisión de la literatura científica existente sobre un tema \item en el trabajo: se revisaron 30 artículos sobre regresión simbólica, búsqueda de raíces y álgebra computacional \item permite identificar qué se ha hecho, qué falta y cómo se posiciona el trabajo
  \end{itemize}

\end{itemize}

\end{document}
